% Drafted against audit/outline.md §9. No new evidence and no new citations.
% Every figure restated here appears with its verification in §6 or §7.

\section{Conclusion}\label{sec:conclusion}

Goldman Sachs discloses one number in Item 7A that depends on something the
filing never publishes. The four category VaRs are levels; the diversification
effect is a relationship, and it is the only line in the table that a reader
outside the firm can test without knowing a single position. This paper tested
it, over ten consecutive fiscal years, with a historical-simulation engine run
at the bank's own disclosed settings on seventeen public risk factors.

The level transfers. A book sized once on 2019 and then left alone averages
41.6\% of summed category VaR against a disclosed 42.0\%, a difference no test
on ten annual observations distinguishes from zero. That corroborates
\citet{perignon2010diversification}, who found disclosed diversification
effects close to independently estimated ones on four US banks and quarterly
data ending in early 2007. Finding it again on a different bank, a different
decade, a different frequency and a different method is worth reporting
precisely because it is not new: a result that survives that much substitution
is a result about disclosure practice rather than about one estimation.

The steadiness does not transfer, and that is what this paper adds. Goldman's
average daily ratio has a year-to-year standard deviation of 2.08 percentage
points across a decade containing a pandemic, a zero-rate period and more than
500 basis points of policy tightening. The proxy's is 6.33, a difference a
variance-ratio test puts at $p = 0.0028$, and the two move oppositely over the
sample: Goldman's second-half mean is 0.8 points above its first, the proxy's
8.5 below.

The firm's second disclosed measure says what kind of failure it is. Goldman
also publishes a period-end snapshot of the same ratio, and that series has a
standard deviation of 6.72 points. Its diversification effect moves a great
deal, and what is stable is the annual mean rather than the quantity. Measured
period-end against period-end, the proxy's variation cannot be distinguished
from Goldman's. What distinguishes them is persistence: the lag-1
autocorrelation of consecutive annual values is $-0.47$ for Goldman's average
daily series and about zero for its period-end series, against roughly $+0.9$
for both proxy series. Goldman's movements do not carry forward; the proxy's
do. The proxy reproduces the level, and its variation is the same size as the
disclosed variation; what it fails to reproduce is the mean reversion, and the
mean reversion is what the disclosed steadiness consists of.

What makes the result hard to dismiss is that it survives being handed the
bank's own method. The filings disclose a five-year scenario window and
age-weighted scenarios, and both were implemented. The window moves the gap the
right way and runs out of room, since five years is what is disclosed and there
is no longer window to appeal to, and the weighting moves it the wrong way,
raising the standard deviation to between 9.46 and 10.04 points while
simultaneously repairing the engine's coverage, from a 6.98\% exception rate
that Kupiec rejects below 0.0001 to 5.45\% at $p = 0.328$. The engine gets
better as a risk model and worse as a reproduction of the disclosed steadiness,
in the same move. Whatever holds Goldman's ratio still is not in the
methodology the filings describe, because that methodology has been applied
here and does not produce it.

Nor is it a matter of having guessed the sleeve directions wrongly. Across all
sixteen sign configurations the two properties trade off: five are at least as
steady as the disclosure and the best of those averages 35.6\%, eight reach
40\% or more and the steadiest of those varies by 3.68 points. None does both.
That trade-off, rather than any single calibration, is what the sign space
establishes. It is a descriptive bound on what this engine can produce rather
than a test --- ten annual observations per configuration will not support one
--- and it is what an outside party can establish from public data.

Taken with the decomposition above, the paper's contribution is two-part: a
replication run at the bank's own disclosed settings, which rules out the
objection that the gap measures a cruder method; and a separation of amplitude
from persistence using a second disclosed series the firm already publishes,
which identifies what the gap actually consists of. Neither requires anything
the filings do not contain.

\subsection*{How much this establishes, and how much it does not}

The claim is narrower than the result might be made to sound. This paper does
not show that Goldman's disclosed diversification effect is implausible,
managed, or anything other than what the firm says it is. It shows that a
static four-sleeve book of public factors cannot reproduce it, under the
method the firm discloses, across the full space of sign configurations
available to such a book. Those are different statements, and the second does
not imply the first.

Three explanations survive the elimination in Section~\ref{sec:assumptions},
and the paper does not rank them because nothing in the evidence ranks them.
The factor count, seventeen against roughly seventy thousand, is the
largest structural difference and its effect on steadiness is simply unknown.
Non-synchronous closes across time zones bias the measurement in the opposite
direction to the factor count, and neither bias is quantified. And exposures
are managed rather than held: in four consecutive filings, FY2021 through
FY2024, Goldman attributes its year-over-year VaR change to volatility moving
one way and exposures moving the other. A book whose notionals were fixed in
2019 has no counterpart to that. It is the one surviving candidate with direct
documentary support, and the only one whose mechanism matches the shape of the
discrepancy: a desk that sheds risk into an unwelcome correlation regime and
restores it afterwards produces a ratio that returns to where it was, which is
precisely what the proxy's does not do.

\subsection*{What would settle it}

Separating those three requires data that is not public. Position-level or
desk-level exposures would do it directly: holding the factor set fixed and
varying only whether positions are repriced or held would isolate exposure
management from the other two, and a richer factor set run against the same
book would isolate the factor count. Supervisory data of the kind
\citet{obrien2017evaluation} worked with would get part of the way, though it
carries the trading outcomes rather than the diversification decomposition this
question needs.

Absent any of that, the bound is what an outside party can establish, and
establishing it is the point. A reader of a 10-K can now say something specific
about the one line in the market-risk table that is testable from outside: its
level is reproducible from public data, the size of its variation is
reproducible too, and its persistence is not, under the bank's own disclosed
window and weighting, swept across the range the filings leave open.
